\chapter{Ejercicios Tema 3}

\definecolor{backcolour}{rgb}{0.95,0.95,0.92}
\definecolor{codegray}{rgb}{0.5,0.5,0.5}
\definecolor{keyblue}{rgb}{0.1,0.1,0.8}


\lstset{
    backgroundcolor=\color{backcolour},   
    basicstyle=\ttfamily\small,
    breakatwhitespace=false,         
    breaklines=true,                 
    captionpos=b,                    
    commentstyle=\color{gray},
    keywordstyle=\color{keyblue},
    numbers=left,                    
    numberstyle=\tiny\color{codegray},
    stepnumber=1,
    numbersep=8pt,                  
    showspaces=false,                
    showstringspaces=false,
    showtabs=false,                  
    tabsize=2,
    frame=single, 
    rulecolor=\color{black!30}
}

\section{Exclusión Mutua de Anillo}

\begin{lstlisting}[language=Pascal, caption=Algoritmo de Exclusión Mutua de Anillo]
    Cliente[1..n] {
        send(solicitar[i]);
        receive(conceder[i]);
        // Seccion critica
        send(liberar[i]);
        ...
    }

    Replica[1..n]{
        ...
        receive(token[i]);
        if (solicitar.size()!=0){
            receive(solicitar[i]);
            send(conceder[i]);
            receive(liberar[i]);
        }
        send(token[(i mod n) + 1]);
    }
\end{lstlisting}
